\SourceLeaf{263}%
insensibili haberi potest. Ita e. g. in triangulo maximo inter ea, quae annis praecedentibus dimensi sumus, puta inter puncta Hohenhagen, Brocken, Inselsberg, ubi excessus summae angulorum fuit \(=14''85348\), calculus sequentes reductiones angulis applicandas prodidit:
\begin{center}
\begin{tabular}{ll}
Hohenhagen \(.\ .\ .\ .\ .\)&\(-4''95113\)\\
Brocken \(.\ .\ .\ .\ .\ .\)&\(-4{,}95104\)\\
Inselsberg\(.\ .\ .\ .\ .\ .\)&\(-4{,}95131\)
\end{tabular}
\end{center}
\Article{29}
\indent Coronidis caussa adhuc comparationem areae trianguli in superficie curva cum area trianguli rectilinei, cuius latera sunt \(a,b,c\), adiiciemus. Aream posteriorem denotabimus per \(\sigma^*\), quae fit \(=\tfrac{1}{2}bc\sin A^*=\tfrac{1}{2}ac\sin B^*=\tfrac{1}{2}ab\sin C^*\)

Habemus, usque ad quantitates ordinis quarti
\[
\sin A^*=\sin A-\tfrac{1}{12}\sigma\cos A\mathord{.}(2\alpha+\beta+\gamma)
\]
sive aeque exacte
\[
\sin A=\sin A^*\mathord{.}(1+\tfrac{1}{24}bc\cos A\mathord{.}(2\alpha+\beta+\gamma))
\]
Substituto hoc valore in formula [9], erit usque ad quantitates sexti ordinis
\[
\begin{aligned}
\sigma=\tfrac{1}{2}bc\sin A^*\mathord{.}\{1&+\tfrac{1}{120}\alpha(3bb+3cc-2bc\cos A)\\
&+\tfrac{1}{120}\beta(3bb+4cc-4bc\cos A)\\
&+\tfrac{1}{120}\gamma(4bb+3cc-4bc\cos A)\}
\end{aligned}
\]
sive aeque exacte
\[
\begin{aligned}
\sigma=\sigma^*\{1&+\tfrac{1}{120}\alpha(aa+2bb+2cc)+\tfrac{1}{120}\beta(2aa+bb+2cc)\\
&+\tfrac{1}{120}\gamma(2aa+2bb+cc)\}
\end{aligned}
\]
Pro superficie sphaerica haec formula sequentem induit formam
\[
\sigma=\sigma^*(1+\tfrac{1}{24}\alpha(aa+bb+cc))
\]
cuius loco etiam sequentem salva eadem praecisione adoptari posse facile confirmatur
\[
\sigma=\sigma^*\BareRoot{\frac{\sin A\mathord{.}\sin B\mathord{.}\sin C}{\sin A^*\mathord{.}\sin B^*\mathord{.}\sin C^*}}
\]
Si eadem formula triangulis in superficie curva non sphaerica applicatur, error generaliter loquendo erit quinti ordinis, sed insensibilis in omnibus triangulis, qualia in superficie telluris dimetiri licet.
\par\medskip\centerline{\rule{.75in}{.3pt}}
